%=====================================================================
\chapter{What Is Inside a DBMS, and Who Does What}\label{ch:whodoeswhat}
%=====================================================================
\locator{1}{Part I --- What a Database Is}

\begin{outcomes}
By the end of this chapter you will be able to:
\begin{tight}
  \item \textbf{Name} the major components of a DBMS and \textbf{state} what
        each one is responsible for.
  \item \textbf{Trace} a single \sqlkw{SELECT} statement through those
        components, in order, from text to answer.
  \item \textbf{Explain} what the data dictionary is, and \textbf{query} it.
  \item \textbf{Distinguish} the four database roles --- administrator,
        designer, application programmer, end user --- by what each one
        decides.
  \item \textbf{Identify} which components SQLite has, which it lacks, and
        which role has to cover for the missing ones.
\end{tight}
\end{outcomes}

\begin{prereq}
\chref{architecture}: the three levels and the two mappings. This chapter is
largely about \emph{who performs} those mappings.
\end{prereq}

\begin{keyterms}
query processor $\bullet$ parser $\bullet$ query optimiser $\bullet$ execution
engine $\bullet$ storage manager $\bullet$ buffer manager $\bullet$ file
manager $\bullet$ transaction manager $\bullet$ lock manager $\bullet$ recovery
manager $\bullet$ authorisation manager $\bullet$ data dictionary $\bullet$
system catalog $\bullet$ metadata $\bullet$ DBA $\bullet$ database designer
$\bullet$ application programmer $\bullet$ end user
\end{keyterms}

\section{A DBMS Is Not One Program}

You type eleven words:

\begin{lstlisting}[language=sql]
SELECT name FROM student WHERE programme = 'BSIT' ORDER BY name;
\end{lstlisting}

and rows come back. Between those two events, six or seven distinct pieces of
software run, each with one job. It is worth knowing them by name, because for
the rest of the course every topic is really a chapter about one of them:
\chref{indexes} is about the optimiser and the storage manager arguing,
\chref{transactions} is about the transaction and recovery managers, and
\chref{project} is about the one component SQLite does not have at all.

\begin{definitionbox}[The components of a DBMS]
The \term{query processor} turns your statement into an answer. It contains a
\term{parser} (is this valid SQL, and what does it refer to?), a
\term{query optimiser} (of the many ways to compute this, which is cheapest?)
and an \term{execution engine} (do it).

\smallskip
The \term{storage manager} stands between the query processor and the disk. It
contains a \term{file manager} (where on disk is this?) and a
\term{buffer manager} (which pages are in memory, and which do I evict?).

\smallskip
The \term{transaction manager} makes groups of operations behave as one, with
a \term{lock manager} to keep concurrent users from corrupting each other's
work and a \term{recovery manager} to put the database back together after a
crash.

\smallskip
The \term{authorisation manager} decides whether this user is allowed to do
this at all.

\smallskip
The \term{data dictionary} --- also called the \term{system catalog} --- is
where the DBMS keeps the schema itself. It is a database about the database.
\end{definitionbox}

\dsfig{%
\begin{tikzpicture}[
  comp/.style={rounded corners=3pt, draw=#1, line width=1pt, fill=#1!7,
               text=#1!60!black, font=\scriptsize\bfseries, align=center,
               text width=2.7cm, minimum height=0.8cm},
  grp/.style={rounded corners=4pt, draw=#1!40, line width=0.8pt, dashed,
              inner sep=6pt},
  ttl/.style={font=\scriptsize\bfseries, text=#1, anchor=west},
  arr/.style={-{Stealth[length=2mm]}, primary!55, line width=0.9pt},
  thin/.style={-{Stealth[length=1.6mm]}, neutral!70, line width=0.6pt}
]
  % users
  \node[comp=secondary] (u) at (0,5.0) {SQL from a user\\or a program};

  % query processor
  \node[comp=primary] (p)  at (0,3.7) {Parser};
  \node[comp=primary] (o)  at (0,2.7) {Optimiser};
  \node[comp=primary] (x)  at (0,1.7) {Execution engine};
  \node[grp=primary, fit=(p)(o)(x)] (qp) {};
  \node[ttl=primary] at ([xshift=2pt]qp.north east) {Query processor};

  % storage manager
  \node[comp=purplehl] (b) at (0,0.3) {Buffer manager};
  \node[comp=purplehl] (f) at (0,-0.7) {File manager};
  \node[grp=purplehl, fit=(b)(f)] (sm) {};
  \node[ttl=purplehl] at ([xshift=2pt]sm.north east) {Storage manager};

  \node[comp=neutral] (disk) at (0,-2.1) {the database file};

  % side services
  \node[comp=tealhl]  (t) at (-4.3,2.2) {Transaction\\manager};
  \node[comp=tealhl]  (l) at (-4.3,0.9) {Lock manager};
  \node[comp=tealhl]  (r) at (-4.3,-0.4) {Recovery\\manager};
  \node[comp=accent]  (a) at (-4.3,3.9) {Authorisation\\manager};
  \node[comp=goldhl]  (d) at (4.3,1.7) {Data dictionary\\{\scriptsize\mdseries
        the schema, stored}};

  \draw[arr] (u) -- (p);
  \draw[arr] (p) -- (o);
  \draw[arr] (o) -- (x);
  \draw[arr] (x) -- (b);
  \draw[arr] (b) -- (f);
  \draw[arr] (f) -- (disk);

  \draw[thin] (a) -- (p);
  \draw[thin] (t) -- (x);
  \draw[thin] (l) -- (b);
  \draw[thin] (r) -- (f);
  \draw[thin, {Stealth[length=1.6mm]}-{Stealth[length=1.6mm]}] (d) -- (o);
  \draw[thin, {Stealth[length=1.6mm]}-{Stealth[length=1.6mm]}] (d) -- (p);
\end{tikzpicture}}%
{The components of a DBMS, and the path a statement takes down through them.
The optimiser and the parser both consult the data dictionary --- the parser
to check that \texttt{student} exists and has a \texttt{name}, the optimiser to
find out how big it is and what indexes it has.}{dbmscomponents}

\section{Following One Statement Down}

\begin{worked}[Eleven words, six components]
\texttt{SELECT name FROM student WHERE programme = 'BSIT' ORDER BY name;}

\begin{tightnum}
  \item \textbf{Authorisation manager.} May this user read
        \texttt{student}? In a server DBMS this is checked first and can stop
        everything here. \emph{In SQLite there is no such check} ---
        if you can open the file, you can read all of it.
  \item \textbf{Parser.} Is the syntax legal? Does \texttt{student} exist?
        Does it have columns \texttt{name} and \texttt{programme}? Is
        \texttt{'BSIT'} comparable with \texttt{programme}? The answers come
        from the \textbf{data dictionary}, not from the data.
  \item \textbf{Optimiser.} Two plans at least. Read all 4{,}000 rows and
        discard the non-BSIT ones, then sort; or, if an index on
        \texttt{programme} exists, use it to fetch only the BSIT rows. And if
        an index on \texttt{name} exists, the \sqlkw{ORDER BY} may be free ---
        indexes are already in order. The optimiser picks by \emph{estimated}
        cost, using statistics from the dictionary. \chref{indexes} measures
        how good those estimates are.
  \item \textbf{Execution engine.} Runs the chosen plan, asking the storage
        manager for pages as it goes.
  \item \textbf{Buffer manager.} Page 7 is already in memory --- no disk read.
        Page 112 is not, so fetch it, and evict something to make room. This
        is the single biggest determinant of how fast the query is, and you
        never see it.
  \item \textbf{File manager.} Page 112 is at byte offset $111 \times 4096$ of
        \texttt{university.db}. Read 4{,}096 bytes.
\end{tightnum}

\smallskip
\textbf{The point of the exercise.} You wrote \emph{what} you wanted. Nothing in
your statement said how. Steps 3 to 6 were chosen for you, and the only way to
influence them is to change the schema --- which is why design, and not
cleverness in query writing, is most of this course.
\end{worked}

\section{The Data Dictionary Is Just More Tables}

The claim that a DBMS stores the schema in the database is easier to believe
once you have queried it.

\begin{examplebox}[The dictionary, in SQLite]
Every SQLite database has a table called \texttt{sqlite\_master}, and you can
select from it like any other:

\begin{lstlisting}[language=sql]
SELECT type, name, tbl_name FROM sqlite_master;
\end{lstlisting}

\begin{lstlisting}
type   name          tbl_name
-----  ------------  ----------
table  student       student
index  idx_name      student
view   fees_view     student
\end{lstlisting}

The \texttt{sql} column holds the exact text of the \sqlkw{CREATE} statement
that made each object --- which is why \texttt{.schema} is instant: it is a
\sqlkw{SELECT}, not a reconstruction.

\smallskip
Other systems do the same thing under other names:
\texttt{information\_schema} in MySQL and PostgreSQL,
\texttt{ALL\_TABLES} in Oracle. The idea is universal; the spelling is not.
\end{examplebox}

\begin{conceptbox}[Why this is more than a curiosity]
Because the schema is data, every tool that works on data works on the schema.
Three consequences you will use:

\begin{tight}
  \item \textbf{A program can discover a database it has never seen.} Anything
        that shows you a list of tables --- a GUI browser, an admin page, a
        reporting tool --- is selecting from the dictionary.
  \item \textbf{The schema can be backed up with the data}, because it is in
        the same file.
  \item \textbf{A dictionary query can audit a schema.} In \chref{constraints}
        you will use one to find every table that declares a foreign key, and
        the point of that exercise is that the declarations are all there and
        none of them are being enforced.
\end{tight}
\end{conceptbox}

\section{The Four Roles}

The components are software. These are people, and the distinction that
matters is not job title --- in a small organisation one person is all four ---
but \emph{which decisions belong to which hat}.

\begin{definitionbox}[Who decides what]
The \term{database administrator} (DBA) owns the running system: installation,
users and privileges, backups, recovery, performance, storage. The DBA decides
\emph{how it runs}.

\smallskip
The \term{database designer} owns the schema: entities, relationships,
constraints, normalisation. The designer decides \emph{what is true}.

\smallskip
The \term{application programmer} owns the programs that use the database.
They decide \emph{how it is reached}.

\smallskip
The \term{end user} owns the questions. They decide \emph{what it is for}, and
they are the only role in the list who need not know that any of the others
exist.
\end{definitionbox}

\begin{examplebox}[One request, four reactions]
\emph{``Students should be able to see their own marks on the portal, but not
anybody else's.''}

\begin{tabular}{@{}L{3.1cm}L{10.9cm}@{}}
\textbf{End user} & States the requirement, including the half of it that is a
prohibition. Notices first when it is wrong. \\
\textbf{Designer} & Is there a reliable way to tell whose marks these are?
Does \relation{enrolment}{\relkey{roll\_no}, \relkey{section\_id}, grade} tie
every mark to exactly one student? If not, the requirement cannot be met at
any price. \\
\textbf{DBA} & Creates the portal's database user and grants it the minimum:
read on a view, and nothing on the base table. \\
\textbf{Programmer} & Writes the query with the logged-in roll number as a
\emph{parameter}, never as pasted-in text --- the difference between those
two is the whole of SQL injection. \\
\end{tabular}

\smallskip
\textbf{Where it breaks.} If the designer gets it wrong, nobody downstream can
fix it: a mark not tied to one student cannot be shown to one student. This is
why design comes before everything else in this course, and why
\chref{ermodel} through \chref{constraints} are four chapters in which you
write almost no SQL.
\end{examplebox}

\begin{alertbox}[In this course you are three of the four, because SQLite has no DBA]
SQLite has no users, no privileges, no server to administer and no process to
tune. Several things on the DBA's list simply do not exist, and the rest move
onto somebody else's:

\begin{tight}
  \item \textbf{Users and privileges:} nobody's. The capability is absent. File
        permissions are the only access control there is, and
        \chref{project} is about what you get back when you move to a server.
  \item \textbf{Backup:} yours, and it is \texttt{copy university.db} --- which
        is genuinely all it takes, and is \chref{transactions}'s subject once
        you ask what happens if you copy it mid-write.
  \item \textbf{Performance:} yours. \texttt{ANALYZE} and
        \texttt{EXPLAIN QUERY PLAN} are the whole toolkit, and
        \chref{indexes} uses both.
  \item \textbf{Schema design, and writing the programs:} yours, which is the
        point --- you will hold the designer's hat and the programmer's hat at
        the same time, and the labs are arranged so you can feel the conflict.
\end{tight}

\smallskip
So: when this handout says ``the DBA would'', read it as ``on a real system
somebody else would, and this term it is you, with one hand tied''.
\end{alertbox}

\begin{pitfall}
\begin{tight}
  \item \textbf{Thinking the optimiser can be instructed.} In SQL you say what,
        never how. If the plan is bad, change the schema or the statistics ---
        not your tone of voice.
  \item \textbf{Confusing the data dictionary with the data.} The dictionary
        holds the schema. \texttt{SELECT COUNT(*) FROM sqlite\_master} counts
        tables, not students.
  \item \textbf{Assuming the roles are four people.} They are four sets of
        decisions. In this course they are all you.
  \item \textbf{Believing ``the application will check it''} is equivalent to a
        constraint in the schema. It is not, and \chref{constraints} is eight
        pages on why.
  \item \textbf{Expecting an authorisation manager in SQLite.} There is none.
        Open the file and you have everything.
\end{tight}
\end{pitfall}

\begin{lab}[Query the dictionary, and watch the optimiser change its mind]
\begin{lstlisting}[language=sql]
-- --- 1. Build something for the optimiser to think about -------
sqlite3 inside.db
CREATE TABLE student (
  roll_no   TEXT PRIMARY KEY,
  name      TEXT NOT NULL,
  programme TEXT NOT NULL
);
-- 300 rows, generated by a recursive CTE (explained in Ch 15):
INSERT INTO student
WITH RECURSIVE n(i) AS (SELECT 1 UNION ALL SELECT i+1 FROM n WHERE i<300)
SELECT printf('2021-IT-%03d', i),
       printf('Student %03d', i),
       CASE i%3 WHEN 0 THEN 'BSIT' WHEN 1 THEN 'BSCS' ELSE 'BSSE' END
FROM n;

-- --- 2. Read the dictionary ------------------------------------
SELECT type, name FROM sqlite_master;
SELECT sql FROM sqlite_master WHERE name = 'student';

-- --- 3. Ask the optimiser what it intends to do ----------------
EXPLAIN QUERY PLAN
  SELECT name FROM student WHERE programme = 'BSIT' ORDER BY name;
-- Write the plan down exactly. Note the words SCAN and ORDER BY.

-- --- 4. Give it an index and ask again -------------------------
CREATE INDEX idx_prog ON student(programme);
EXPLAIN QUERY PLAN
  SELECT name FROM student WHERE programme = 'BSIT' ORDER BY name;
-- The plan changed and you did not touch the query.

-- --- 5. Make the sort free too ---------------------------------
CREATE INDEX idx_name ON student(name);
EXPLAIN QUERY PLAN SELECT name FROM student ORDER BY name;
-- Does "USE TEMP B-TREE FOR ORDER BY" still appear?

-- --- 6. Look at what you have added to the dictionary ----------
SELECT type, name, tbl_name FROM sqlite_master;
-- Three new rows you never inserted.
\end{lstlisting}

\smallskip
\textbf{What to notice.} In steps 3 to 5 the statement was identical and the
plan was not. Everything that changed, changed in the schema. Keep the three
plans --- \chref{indexes} asks you to compare them against measured times, and
one of them will turn out to be a worse idea than it looks.
\end{lab}

\begin{checkpoint}
\begin{tightnum}
  \item A query runs in 40 milliseconds on Monday and 4 seconds on Friday, on
        the same data and with no change to the statement or the schema. Name
        two components that could be responsible and say how you would tell
        them apart.
  \item The data dictionary is stored as tables in the database it describes.
        Give one real advantage of that and one problem it creates.
  \item For each of the four roles, state one decision that is theirs alone,
        and then say what happens to that decision in this course --- where
        SQLite provides no DBA at all.
\end{tightnum}
\end{checkpoint}

\begin{chaptersummary}
\begin{tight}
  \item A DBMS is several cooperating components: a \textbf{query processor}
        (parser, optimiser, execution engine), a \textbf{storage manager}
        (buffer manager, file manager), a \textbf{transaction manager} with
        lock and recovery managers, an \textbf{authorisation manager}, and the
        \textbf{data dictionary}.
  \item A statement travels down through them. You specify \emph{what}; steps
        3 onward --- plan, pages, bytes --- are chosen for you.
  \item The parser consults the dictionary to check names and types; the
        optimiser consults it for sizes and indexes.
  \item The \textbf{data dictionary} is ordinary tables. In SQLite it is
        \texttt{sqlite\_master}, and \texttt{.schema} is a \sqlkw{SELECT} over
        it.
  \item Because the schema is data, programs can discover a database they have
        never seen, and a query can audit a schema.
  \item Four roles, defined by their decisions: the \textbf{DBA} decides how it
        runs, the \textbf{designer} decides what is true, the
        \textbf{programmer} decides how it is reached, the \textbf{end user}
        decides what it is for.
  \item A designer's error cannot be repaired downstream, which is why design
        precedes SQL in this course.
  \item SQLite has no authorisation manager and no DBA role. Backup,
        performance and design all land on you, and \chref{project} is the
        chapter about what a server gives back.
\end{tight}
\end{chaptersummary}

\begin{reviewq}
  \item \textbf{(MCQ)} Which component decides whether to use an index?
        (a)~the parser; (b)~the buffer manager; (c)~the query optimiser;
        (d)~the application programmer.
  \item \textbf{(MCQ)} \texttt{sqlite\_master} contains: (a)~the rows of every
        table; (b)~the schema of the database; (c)~the query log;
        (d)~the list of users.
  \item \textbf{(Short)} Explain why the parser needs the data dictionary, with
        a specific example of a statement it would reject because of
        something it found there.
  \item \textbf{(Short)} What does the buffer manager do, and why can it make
        the same query take very different times on identical data?
  \item \textbf{(Short)} Which of the four roles is responsible when a database
        permits two students to have the same roll number? Justify your answer
        in two sentences.
  \item \textbf{(Applied)} Take the statement \texttt{SELECT s.name, e.grade
        FROM student s JOIN enrolment e ON s.roll\_no = e.roll\_no WHERE
        e.grade = 'A';} and walk it through all six components, saying at each
        step what information that component needs and where it gets it.
  \item \textbf{(Applied)} A college wants one database used by the
        examinations office, the fees office and 4{,}000 students. Divide the
        work among the four roles: list what each must do before the system
        opens, then mark each item that SQLite cannot support and say what you
        would do about it.
\end{reviewq}
